Proponer una estrategia de extincion y ver si ese es un equilibrio

proponer una estrategia de threhsold y ver si converge al equilibrio 

meanfield games (es como un fluid game) revisar paper de gabriel W. 

Que el p0 sea endogeno? :

- La estrategia mia debe ser consistente con el comportamiento del buyer y same for el otro seller

Idea del modelo: 

- sellers simetricos 
- escribir general con el rho quizas y cuando rho es 1 entonces no hay forgetting 

- Argumentos para justificar que estamos usando un modelo fluido 
- investigar si es que en el caso estocastico podemos hacer una transformacion parecida a lo del F, cosa que dependa solo del estado 

Tareas:

1. definir el modelo de competencia
2. Argumentos para justificar el modelo fluido y cerrar las pruebas (documento con las proposiciones y pruebas del modelo fluido completo)
3. Investigar si es que en el caso estocastico podemos hacer una transformacion parecida a lo del F, cosa que dependa solo del estado 


-------------------------------------------

- Ver cuanta es la diferencia de valores cuando usamos la politica del modelo fluido vs la politica del problema estocastico
- Aclarar por que con el caso base no se puede resolver y por que la dimension es mas compleja que en el caso sin juego (tambien responder a si hay convergencia en el equilibrio)
- Tener una buena justificación de por qué miramos el oblivious equilibrium 
- Estudiar este modelo de competencia en fluido 

Next steps concretos: 

1. simular el juego estocastico de oblivious equilibrium
2. que cosas nos gustaria probar (existe un equilibrio?, quizas resultados similares al caso del monopolista)
3. Como se comporta este juego cuando tienes dos retailers y los dos se comportan como monopolistas 
    - asume que estoy compitiendo con una calidad fija p0 
    - p0 y p0' deberian ser consistentes con la estragia del otro tipo
    - la estrategia del primer proveedor depende de un p0 pero en realidad p0 es una consecuencia del juego y se pide consistencia 
    - probar como calibrar este p0 

------------------------------------------------
19 de octubre

-simular ahora suponiendo que el competidor tiene una distribucion condicional en tu propio estado
- plotear la estrategia de del monopolista cuando usamos el p0 de equilibrio del OEp0 
- buscar como comparan en le paper de gabriel el MPe vs el oblvious 


------------------------------------------------------------

Resolver el MPE considerando que te quedas con los ultimos N resultados y armar un LP. 
quizas comparar el price of anarchy. 

definicion del modelo de competencia
Se puede escribir un modelo fluido para competencia ? 

